\subsection*{Event Basics}

The meaning of an event varies between programming languages. In Croqtile, an \textbf{event} coordinates asynchronous code blocks at different execution levels.

A Croqtile event has two states, \textbf{SET} and \textbf{UNSET}. A newly declared event starts in the UNSET state. For example:

\begin{lstlisting}[language=croqtile,caption={Declaring a shared event}]
shared event e;
\end{lstlisting}

A \texttt{shared} event is declared inside a \texttt{parallel-by} block and used by its parallel instances. Events can also use \texttt{global} or \texttt{local} storage. The target platform determines the constraints on their declarations and references.

In addition, you may declare an event array following C-style syntax:

\begin{lstlisting}[language=croqtile,caption={Declaring an event array}]
shared event e[4], e1;
\end{lstlisting}

A comma-separated expression can declare both an event and an event array on one line.

The \texttt{trigger} statement changes an event from \textbf{UNSET} to \textbf{SET}:

\begin{lstlisting}[language=croqtile,caption={Triggering an event array}]
shared event e[4], e1;
trigger e, e1;
\end{lstlisting}

The comma-separated statement triggers every event in array \texttt{e} and the single event \texttt{e1}.

The \texttt{wait} statement blocks the current instance while the event is \textbf{UNSET}. Execution resumes once the event becomes \textbf{SET}.

Triggering an event is what unblocks the code waiting on a \emph{wait statement}. For example,

\begin{lstlisting}[language=croqtile,caption={Waiting on an event}]
parallel p by 2 {
  shared event e;                          // init 'e' to 'UNSET'
  inthreads.async (p == 1) { wait e; }     // wait-thread
                                           // once unblocked, change 'e' to 'UNSET'
  inthreads.async (p == 0) { trigger e; }  // trigger-thread
                                           // convert 'e' to 'SET'
  sync.barrier : thread;                   // sync point
}
\end{lstlisting}

In this example, the triggering instance sets \texttt{e}, allowing the waiting instance to continue. Both then reach \texttt{sync.barrier}. The event establishes a dependency between the two asynchronous regions: the trigger must occur before the wait can finish.

In Croqtile semantics, \textbf{an event is always auto-reset by its \texttt{wait} statement}. That means when a \texttt{wait} statement is unblocked, the events that are waited are changed to \emph{UNSET} state.
